% Vietnamese translation for OI Public Library
% Source-PDF: IOI中国国家候选队论文/2008/Day2/6.余林韵《运用化归思想解决信息学中的数列问题》/附件数学相关知识.pdf
\documentclass[11pt]{article}
\usepackage{oipl-vn}

\title{Kiến thức toán học liên quan}
\author{Dư Lâm Vận}
\date{22 tháng 1 năm 2008}

\begin{document}
\maketitle

\origtitle{Mathematical related knowledge}
\sourcepath{IOI-China-Candidate-Team-Papers/2008/Day2/6-Yu-Linyun/math-related-knowledge.pdf}

\section{Định nghĩa dãy số}

Một dãy các số được sắp xếp theo một thứ tự nhất định được gọi là một
\term{dãy số}. Thật ra, nhìn từ quan điểm hàm số, đối với một hàm có miền xác
định là tập số nguyên dương $\mathbb{N}^{*}$ hoặc một tập con hữu hạn
$\{1,2,\ldots,n\}$ của nó, dãy số chính là dãy các giá trị hàm tương ứng khi
biến độc lập lần lượt nhận các giá trị từ nhỏ đến lớn.

\section{Số hạng tổng quát và truy hồi}

Công thức số hạng tổng quát và công thức truy hồi là hai phương pháp quan
trọng để cho trước một dãy số. Trong các bài toán tin học, dãy số thường được
cho bằng công thức truy hồi, hoặc được giải bằng cách sử dụng công thức truy
hồi.

\subsection{Công thức số hạng tổng quát}

Nếu quan hệ giữa số hạng thứ $n$ của dãy $\{a_n\}$ và $n$ có thể được biểu
diễn bằng một công thức, thì công thức ấy được gọi là \term{công thức số hạng
tổng quát} của dãy số đó.

\subsection{Dãy truy hồi}

\subsubsection{Dãy truy hồi}

Quan hệ giữa các số hạng liên tiếp của một dãy số được gọi là quan hệ truy
hồi. Dãy số được xác định bởi quan hệ truy hồi được gọi là dãy truy hồi.

\subsubsection{Dãy truy hồi bậc \texorpdfstring{$k$}{k}}

Dãy số $\{a_n\}$ được xác định bởi các giá trị ban đầu và phương trình
\eqref{eq:kth-recurrence} được gọi là dãy truy hồi bậc $k$.
\begin{equation}
  a_{n+k}=F(a_{n+k-1},\ldots,a_n).
  \label{eq:kth-recurrence}
\end{equation}

Đặc biệt, khi \eqref{eq:kth-recurrence} có dạng
\eqref{eq:linear-recurrence}, dãy số $\{a_n\}$ được gọi là dãy truy hồi tuyến
tính bậc $k$ với hệ số hằng. Ở đây $c_1,c_2,\ldots,c_k$ là các hằng số và
$c_k\ne 0$. Nếu hàm $f(n)=0$, thì dãy số $\{a_n\}$ được xác định bởi
\eqref{eq:linear-recurrence} được gọi là dãy truy hồi tuyến tính thuần nhất
bậc $k$ với hệ số hằng.
\begin{equation}
  a_{n+k}=c_1a_{n+k-1}+c_2a_{n+k-2}+\cdots+c_ka_n+f(n).
  \label{eq:linear-recurrence}
\end{equation}

\subsection{Cấp số nhân}

Dãy số thỏa mãn công thức truy hồi $a_{n+1}=qa_n$, trong đó $q$ là hằng số
khác $0$, được gọi là \term{cấp số nhân}. Hằng số $q$ được gọi là công bội của
cấp số nhân $(q\ne 0)$.

Vì trong một cấp số nhân $\{a_n\}$, kể từ số hạng thứ hai, tỉ số giữa mỗi số
hạng và số hạng đứng ngay trước nó đều bằng công bội $q$, nên mỗi số hạng đều
bằng số hạng trước nhân với $q$. Do đó:
\[
  a_2=a_1q,
\]
\[
  a_3=a_2q=(a_1q)q=a_1q^2,
\]
\[
  a_4=a_3q=(a_1q^2)q=a_1q^3,
\]
\[
  \cdots
\]
Từ đó suy ra
\[
  a_n=a_1q^{n-1}.
\]

Tổng $n$ số hạng đầu của một cấp số nhân có công bội khác $1$: đặt
$S_n=a_1+a_2+a_3+\cdots+a_n$, khi đó:
\[
  S_n=a_1+a_1q+a_1q^2+\cdots+a_1q^{n-1},
\]
\[
  qS_n=0+a_1q+a_1q^2+\cdots+a_1q^{n-1}+a_1q^n.
\]
Vì vậy:
\[
  (1-q)S_n=a_1-a_1q^n,
\]
\[
  S_n=\frac{a_1(1-q^n)}{1-q}.
\]

Dãy số thỏa mãn công thức truy hồi $a_{n+2}=2a_{n+1}-a_n$ được gọi là
\term{cấp số cộng}. Nó và cấp số nhân là những dãy truy hồi đơn giản nhất.

\section{Dãy tuần hoàn}

\subsection{Khái niệm dãy tuần hoàn}

\begin{enumerate}[label=\alph*)]
  \item Đối với dãy số $a_n$, nếu tồn tại các số tự nhiên xác định $T$ và
  $n_0$ sao cho với mọi $n\ge n_0$ luôn có $a_{n+T}=a_n$, thì $a_n$ được gọi
  là dãy tuần hoàn kể từ số hạng thứ $n_0$, với chu kỳ $T$. Khi $n_0=1$, ta
  gọi $a_n$ là dãy thuần tuần hoàn; khi $n_0\ge 2$, ta gọi $a_n$ là dãy hỗn
  tuần hoàn.

  \item Giả sử $\{a_n\}$ là một dãy số nguyên, $m$ là một số tự nhiên đã chọn
  lớn hơn $1$. Nếu $b_n$ là phần dư khi chia $a_n$ cho $m$, tức
  \[
    b_n\equiv a_n \pmod{m},\qquad
    b_n\in\{0,1,2,\ldots,m-1\},
  \]
  thì dãy $\{b_n\}$ được gọi là dãy modulo $m$ của $\{a_n\}$, ký hiệu là
  $a_n \pmod{m}$.
\end{enumerate}

Nếu dãy modulo $\{a_n \pmod{m}\}$ có tính tuần hoàn, thì $\{a_n\}$ được gọi là
dãy tuần hoàn theo modulo $m$.

\subsection{Các tính chất và kết luận quan trọng về dãy tuần hoàn}

\begin{enumerate}[label=\alph*)]
  \item Dãy tuần hoàn là dãy vô hạn, và miền giá trị của nó là một tập hữu hạn.

  \item Nếu $T$ là một chu kỳ của $\{a_n\}$, thì với mọi $k\in\mathbb{N}^{*}$,
  $kT$ cũng là một chu kỳ của $\{a_n\}$.

  \item Dãy tuần hoàn nhất định có chu kỳ dương nhỏ nhất.

  \item Nếu $T$ là chu kỳ dương nhỏ nhất của dãy tuần hoàn $\{a_n\}$, còn
  $T'$ là một chu kỳ bất kỳ của $\{a_n\}$, thì $T\mid T'$.

  \item Một dãy truy hồi tuyến tính thuần nhất bậc $k$ với hệ số hằng bất kỳ
  thỏa mãn quan hệ truy hồi tuyến tính
  \[
    a_{n+k}=c_1a_{n+k-1}+c_2a_{n+k-2}+\cdots+c_ka_n+f(n)
  \]
  là một dãy tuần hoàn theo modulo. Đặc biệt, dãy Fibonacci $\{a_n\}$:
  \[
    \{a_0=a_1=1,\ a_{n+1}=a_n+a_{n-1}\ (n\ge 1)\},
  \]
  với mọi số tự nhiên $p$, $\{a_n \pmod{p}\}$ là một dãy thuần tuần hoàn.
\end{enumerate}

\section{Quy nạp toán học}

\subsection{Định lý quy nạp toán học}

Nếu một mệnh đề $P(n)$ về số tự nhiên $n$ thỏa mãn các điều kiện sau:
\begin{enumerate}
  \item $P(1)$ đúng.
  \item Nếu $P(k)$ đúng, thì $P(k+1)$ đúng.
\end{enumerate}
thì $P(n)$ đúng với mọi $n\in\mathbb{N}^{+}$.

Từ định lý trên, ta thấy các bước dùng quy nạp toán học để chứng minh một
mệnh đề liên quan đến số tự nhiên là:
\begin{enumerate}
  \item Chứng minh kết luận đúng khi $n$ nhận giá trị đầu tiên $n_0$, chẳng
  hạn $n_0=1$ hoặc $n_0=2$.

  \item Giả sử kết luận đúng khi $n=k$ $(k\in\mathbb{N}^{+}$ và $k\ge n_0)$,
  rồi chứng minh kết luận cũng đúng khi $n=k+1$.
\end{enumerate}
Sau khi hoàn thành hai bước này, có thể khẳng định mệnh đề đúng với mọi số tự
nhiên $n$ bắt đầu từ $n_0$.

\subsection{Quy nạp toán học thứ hai}

Nếu một mệnh đề $P(n)$ về số tự nhiên $n$ thỏa mãn các điều kiện sau:
\begin{enumerate}
  \item $P(1)$ đúng.
  \item Với $k\in\mathbb{N}^{+}$, nếu $P(n)$ đúng với mọi $n\le k$, thì
  $P(k+1)$ đúng.
\end{enumerate}
thì $P(n)$ đúng với mọi $n\in\mathbb{N}^{+}$.

\subsection{Quy nạp ngược}

Quy nạp ngược còn được gọi là quy nạp lùi, do nhà toán học Pháp Cauchy sử
dụng đầu tiên. Phương pháp quy nạp ngược được phát biểu như sau: giả sử
$P(n)$ là một mệnh đề về số tự nhiên $n\in\mathbb{N}^{+}$. Nếu:
\begin{enumerate}
  \item $P(n)$ đúng với vô hạn nhiều số tự nhiên $n$.
  \item Giả sử $P(k+1)$ đúng, ta có thể suy ra $P(k)$ đúng.
\end{enumerate}
thì mệnh đề $P(n)$ đúng với mọi $n\in\mathbb{N}^{+}$.

\subsection{Quy nạp bập bênh}

Giả sử $P(n)$ và $Q(n)$ là hai mệnh đề liên quan đến số tự nhiên
$n\in\mathbb{N}^{+}$. Nếu:
\begin{enumerate}
  \item $P(1)$ đúng.
  \item Giả sử $P(k)$ đúng, ta có thể suy ra $Q(k)$ đúng; giả sử $Q(k)$ đúng,
  ta có thể suy ra $P(k+1)$ đúng.
\end{enumerate}
thì với mọi $n\in\mathbb{N}^{+}$, cả $P(n)$ và $Q(n)$ đều đúng.

\subsection{Quy nạp toán học hai chiều}

Khi chứng minh một mệnh đề $P(m,n)$ liên quan đến hai số tự nhiên độc lập, có
thể tiến hành theo hình thức sau:
\begin{enumerate}
  \item Chứng minh $P(1,m)$ đúng với mọi số tự nhiên $m$, và $P(n,1)$ đúng
  với mọi số tự nhiên $n$ $(m,n\in\mathbb{N}^{+})$.

  \item Giả sử $P(n+1,m)$ và $P(n,m+1)$ đúng, từ đó suy ra $P(n+1,m+1)$ đúng.
\end{enumerate}
Khi đó $P(m,n)$ đúng với mọi số nguyên dương $m,n$. Phương pháp này được gọi
là quy nạp toán học hai chiều.

\end{document}
